\chapter{人机交互强化策略分析与系统优化}

\section{风险识别与告警分级策略}

固定翼飞行器着陆阶段的风险并非在接地瞬间才会出现，而是在进近、拉平、接地区域逐渐累积。因此人机交互强化首先需要解决的问题便是及时识别关键风险，再通过分级告警的方式提醒操纵者。从仿真结果可知，复杂地形、下沉率偏大、侧偏增加以及姿态波动加剧都会提高着陆风险，特别是在山脊地形和阶跃地形中，这种风险更为显著。

风险识别应当围绕飞行速度、下沉率、俯仰角、滚转角、侧偏量、航向误差和跑道相对位置等状态量进行。其中，速度主要反映是否存在失速或过快接地的风险，下沉率主要反映接地冲击风险，俯仰角与滚转角主要反映姿态超限风险，侧偏量和航向误差反映跑道对准状态。为提高提示效果，本文将告警划分为三级：第一级是提示级，表示状态偏离目标但尚未接近边界；第二级是警示级，意味着关键参数接近限制范围；第三级是危险级，表明系统已经接近或进入明显不利窗口，此时应触发保护逻辑或复飞机制。

\section{状态提示、操纵辅助与约束控制设计}

在人机交互强化设计中，风险识别仅仅是基础，真正对改善着陆品质起关键作用的是状态提示、操纵辅助和约束控制三者的协同配合。本文将人机交互强化功能归纳为"提示—辅助—约束"三层结构。第一层为状态提示层，第二层是操纵辅助层，第三层是约束控制层。这三层结构按照由弱到强的顺序逐渐介入，既能保留操纵者的参与感，又能在高风险阶段提供必要保护。

状态提示层设计的关键在于追求"简单、直接、可理解"。着陆末段操纵负荷较高，系统不应给操纵者传递过多冗余信息，而应着重展示最关键的状态量。具体包括离地高度、速度、下沉率、俯仰角、滚转角、侧偏量、跑道剩余距离和目标接地点相对位置等。并且可以借助颜色变化、趋势箭头或文字提示等形式，对"偏大""偏小""趋于危险"等状态予以直观呈现。

操纵辅助层的要点是"给予方向，而非进行替代"。当飞行器状态接近风险边界但尚未进入危险状态时，系统应向操纵者输出建议信息，如建议增大或减小俯仰修正、提醒保持中心线、提示优先控制下沉率等。对于末段着陆而言，这种辅助具有重要意义，因为操纵者常在高度快速变化时出现操作犹豫或过度修正。

约束控制层为人机交互强化的最后一道保障。当根据风险识别判定飞行器已逼近危险边界，或操纵输入可能导致姿态快速恶化时，系统需对部分输入实施限幅、整形或强制约束。例如，在近地阶段对副翼和方向舵的过大修正予以限制，防止横向姿态突变；速度过低时提高油门下限，避免进入危险低速区；下沉率严重超限时，触发更强的拉平控制或复飞保护。约束控制的目的并非完全接管操纵，而是在关键窗口防止系统进入不可恢复状态。

\begin{table}[H]
\centering
\caption{人机交互强化三层功能设计}
\begin{tabular}{|l|l|l|}
\hline
\textbf{层级} & \textbf{功能} & \textbf{具体内容} \\
\hline
第一层 & 状态提示 & 离地高度、速度、下沉率、姿态角、侧偏量等核心状态显示 \\
\hline
第二层 & 操纵辅助 & 建议俯仰修正、提醒保持中心线、提示控制下沉率 \\
\hline
第三层 & 约束控制 & 限幅保护、油门下限、复飞触发 \\
\hline
\end{tabular}
\end{table}

\section{关键参数敏感性分析}

为明确哪些参数对着陆品质和交互强化效果影响最为显著，本文针对关键参数进行敏感性分析。结合模型和工况结果，选择目标速度、下沉率阈值、拉平高度、跑道捕获距离、侧偏约束和滚转角约束作为敏感性分析对象。这些参数在控制逻辑中属于关键边界条件，且与交互提示、保护动作的触发时机直接相关。

\begin{table}[H]
\centering
\caption{关键参数敏感性分析}
\begin{tabular}{|l|l|}
\hline
\textbf{参数名称} & \textbf{对系统的影响程度} \\
\hline
目标速度 & 高 \\
\hline
下沉率阈值 & 高 \\
\hline
拉平高度 & 高 \\
\hline
侧偏约束 & 中 \\
\hline
滚转角约束 & 中 \\
\hline
跑道捕获距离 & 中 \\
\hline
\end{tabular}
\end{table}

\section{改进建议与工程可实现性讨论}

借助敏感性分析能够发现，目标速度、下沉率阈值、拉平高度对纵向品质的影响最为显著，而侧偏量和滚转角约束对横向稳定和跑道对准的影响更为突出。这表明后续系统优化不应平均地对所有参数进行调整，而应优先围绕高敏感性参数展开。

根据工况分析和敏感性分析结果，本文认为人机交互强化的系统优化可从四个方面着手。其一，优化告警逻辑，进一步提高告警分级的针对性和可解释性，避免提示过多造成信息干扰或提示过晚失去干预价值。其二，优化状态显示方式，突出离地高度、下沉率、速度和侧偏量等核心状态，让操纵者在末段能够迅速抓住重点。其三，优化操纵辅助策略，在不同着陆阶段设置不同辅助优先级：进近阶段优先给出轨迹与对准提示，拉平阶段优先给出下沉率和俯仰控制提示，接地区域优先给出稳定接地与方向保持提示。其四，优化约束控制边界，对于拉平高度、速度下限、横向约束阈值等高敏感参数，结合工况特征进行分阶段调整，而非采用固定不变的全程阈值。

从工程可实现角度来看，本文提出的人机交互强化策略具备良好的落地基础。首先，速度、高度、姿态和侧偏量等关键状态量均可借助现有飞控系统和导航感知系统获取，无需额外建立复杂硬件模式。其次，告警分级、状态提示和操纵辅助本质上是软件层的逻辑优化，可通过飞控显示界面或地面站界面实现，开发成本相对较低。最后，约束控制和复飞保护虽然对控制逻辑有较高要求，但在已有分阶段控制基础上增加边界判断模块即可实现，不需要重新建立整套控制系统。因此，本文提出的优化策略既有理论分析意义，也具有一定的工程应用前景。

人机交互强化的价值不在于替代飞行控制本身，而在于借助"提前识别—及时提示—适度辅助—必要约束"的方式，提高固定翼飞行器在复杂工况下的着陆稳定性和安全性。
